%! Author = alejandrogonzalezturrubiates
%! Date = 02/10/25

\paragraph{Multiplicación de matrices}

% ===== Ejercicio 1 =====
\begin{ejercicio}
Calcule $A \times B$, donde
\[
A=\begin{bmatrix} 1 & 2 \\ 3 & 4 \end{bmatrix}, \quad
B=\begin{bmatrix} 5 & 6 \\ 7 & 8 \end{bmatrix}.
\]
\end{ejercicio}
\muestraSolucion{
\[
AB=\begin{bmatrix} 1(5)+2(7) & 1(6)+2(8) \\ 3(5)+4(7) & 3(6)+4(8) \end{bmatrix}=
\begin{bmatrix} 19 & 22 \\ 43 & 50 \end{bmatrix}.
\]
}

% ===== Ejercicio 2 =====
\begin{ejercicio}
Calcule $A \times B$, donde
\[
A=\begin{bmatrix} 2 & 0 \\ -1 & 3 \end{bmatrix}, \quad
B=\begin{bmatrix} 4 & -2 \\ 1 & 5 \end{bmatrix}.
\]
\end{ejercicio}
\muestraSolucion{
\[
AB=\begin{bmatrix} 8 & -4 \\ -1 & 17 \end{bmatrix}.
\]
}

% ===== Ejercicio 3 =====
\begin{ejercicio}
Calcule $A \times B$, donde
\[
A=\begin{bmatrix} 3 & -1 \\ 2 & 4 \end{bmatrix}, \quad
B=\begin{bmatrix} 0 & 5 \\ 2 & 1 \end{bmatrix}.
\]
\end{ejercicio}
\muestraSolucion{
\[
AB=\begin{bmatrix} -2 & 14 \\ 8 & 14 \end{bmatrix}.
\]
}

% ===== Ejercicio 4 =====
\begin{ejercicio}
Calcule $A \times B$, donde
\[
A=\begin{bmatrix} 1 & 3 \\ 5 & 7 \end{bmatrix}, \quad
B=\begin{bmatrix} 2 & 4 \\ 6 & 8 \end{bmatrix}.
\]
\end{ejercicio}
\muestraSolucion{
\[
AB=\begin{bmatrix} 20 & 28 \\ 52 & 76 \end{bmatrix}.
\]
}

% ===== Ejercicio 5 =====
\begin{ejercicio}
Calcule $A \times B$, donde
\[
A=\begin{bmatrix} 1 & 2 & 3 \\ 4 & 5 & 6 \end{bmatrix}, \quad
B=\begin{bmatrix} 7 & 8 \\ 9 & 10 \\ 11 & 12 \end{bmatrix}.
\]
\end{ejercicio}
\muestraSolucion{
\[
AB=\begin{bmatrix} 58 & 64 \\ 139 & 154 \end{bmatrix}.
\]
}

% ===== Ejercicio 6 =====
\begin{ejercicio}
Calcule $A \times B$, donde
\[
A=\begin{bmatrix} 2 & 4 \\ 1 & -3 \\ 0 & 5 \end{bmatrix}, \quad
B=\begin{bmatrix} 3 & 1 & 2 \\ 4 & -1 & 0 \end{bmatrix}.
\]
\end{ejercicio}
\muestraSolucion{
\[
AB=\begin{bmatrix} 22 & -2 & 4 \\ -9 & 4 & 2 \\ 20 & -5 & 0 \end{bmatrix}.
\]
}

% ===== Ejercicio 7 =====
\begin{ejercicio}
Calcule $A \times B$, donde
\[
A=\begin{bmatrix} 1 & 0 & 2 \\ -1 & 3 & 1 \end{bmatrix}, \quad
B=\begin{bmatrix} 3 & 1 \\ 2 & 1 \\ 1 & 0 \end{bmatrix}.
\]
\end{ejercicio}
\muestraSolucion{
\[
AB=\begin{bmatrix} 5 & 1 \\ 4 & 2 \end{bmatrix}.
\]
}

% ===== Ejercicio 8 =====
\begin{ejercicio}
Calcule $A \times B$, donde
\[
A=\begin{bmatrix} 2 & -1 & 3 \\ 4 & 0 & 1 \end{bmatrix}, \quad
B=\begin{bmatrix} 1 & 5 \\ 2 & 0 \\ 3 & 1 \end{bmatrix}.
\]
\end{ejercicio}
\muestraSolucion{
\[
AB=\begin{bmatrix} 9 & 13 \\ 7 & 21 \end{bmatrix}.
\]
}

% ===== Ejercicio 9 =====
\begin{ejercicio}
Calcule $A \times B$, donde
\[
A=\begin{bmatrix} 1 & 2 \\ 3 & 4 \\ 5 & 6 \end{bmatrix}, \quad
B=\begin{bmatrix} 7 & 8 & 9 \\ 10 & 11 & 12 \end{bmatrix}.
\]
\end{ejercicio}
\muestraSolucion{
\[
AB=\begin{bmatrix} 27 & 30 & 33 \\ 61 & 68 & 75 \\ 95 & 106 & 117 \end{bmatrix}.
\]
}

% ===== Ejercicio 10 =====
\begin{ejercicio}
Calcule $A \times B$, donde
\[
A=\begin{bmatrix} 2 & 1 & 0 \\ 3 & -1 & 4 \\ 0 & 5 & 2 \end{bmatrix}, \quad
B=\begin{bmatrix} 1 & 2 & 3 \\ 0 & 1 & 4 \\ 5 & 6 & 0 \end{bmatrix}.
\]
\end{ejercicio}
\muestraSolucion{
\[
AB=\begin{bmatrix} 2 & 5 & 10 \\ 23 & 29 & 5 \\ 10 & 17 & 20 \end{bmatrix}.
\]
}

% ===== Ejercicio 11 =====
\begin{ejercicio}
Calcule $A \times B$, donde
\[
A=\begin{bmatrix} 3 & 2 \\ 1 & 0 \\ -2 & 5 \end{bmatrix}, \quad
B=\begin{bmatrix} 1 & -1 & 2 \\ 3 & 0 & 4 \end{bmatrix}.
\]
\end{ejercicio}
\muestraSolucion{
\[
AB=\begin{bmatrix} 9 & -3 & 14 \\ 1 & -1 & 2 \\ 13 & 2 & 16 \end{bmatrix}.
\]
}

% ===== Ejercicio 12 =====
\begin{ejercicio}
Calcule $A \times B$, donde
\[
A=\begin{bmatrix} 4 & -3 \\ 2 & 5 \end{bmatrix}, \quad
B=\begin{bmatrix} 7 & 1 \\ 0 & -2 \end{bmatrix}.
\]
\end{ejercicio}
\muestraSolucion{
\[
AB=\begin{bmatrix} 28 & 10 \\ 14 & -8 \end{bmatrix}.
\]
}

% ===== Ejercicio 13 =====
\begin{ejercicio}
Calcule $A \times B$, donde
\[
A=\begin{bmatrix} 1 & 1 & 1 \\ 0 & 1 & 2 \\ 2 & 1 & 0 \end{bmatrix}, \quad
B=\begin{bmatrix} 2 & 0 & 1 \\ 1 & 2 & 1 \\ 0 & 3 & 2 \end{bmatrix}.
\]
\end{ejercicio}
\muestraSolucion{
\[
AB=\begin{bmatrix} 3 & 5 & 4 \\ 1 & 8 & 5 \\ 5 & 2 & 3 \end{bmatrix}.
\]
}

% ===== Ejercicio 14 =====
\begin{ejercicio}
Calcule $A \times B$, donde
\[
A=\begin{bmatrix} -1 & 2 \\ 4 & 0 \\ 3 & -2 \end{bmatrix}, \quad
B=\begin{bmatrix} 5 & 1 & 0 \\ -3 & 2 & 4 \end{bmatrix}.
\]
\end{ejercicio}
\muestraSolucion{
\[
AB=\begin{bmatrix} -11 & 3 & 8 \\ 20 & 4 & 0 \\ 21 & -1 & -8 \end{bmatrix}.
\]
}

% ===== Ejercicio 15 =====
\begin{ejercicio}
Calcule $A \times B$, donde
\[
A=\begin{bmatrix} 1 & 2 & -1 \\ 0 & 3 & 4 \end{bmatrix}, \quad
B=\begin{bmatrix} 2 & 1 \\ -3 & 0 \\ 1 & 5 \end{bmatrix}.
\]
\end{ejercicio}
\muestraSolucion{
\[
AB=\begin{bmatrix} -5 & -4 \\ -5 & 20 \end{bmatrix}.
\]
}

% ===== Ejercicio 16 =====
\begin{ejercicio}
Calcule $A \times B$, donde
\[
A=\begin{bmatrix} 2 & 4 & 6 \\ 1 & 3 & 5 \\ 0 & 2 & 4 \end{bmatrix}, \quad
B=\begin{bmatrix} 1 & 0 & 2 \\ 3 & 1 & 4 \\ 5 & 2 & 6 \end{bmatrix}.
\]
\end{ejercicio}
\muestraSolucion{
\[
AB=\begin{bmatrix} 44 & 16 & 56 \\ 35 & 13 & 44 \\ 26 & 10 & 32 \end{bmatrix}.
\]
}

% ===== Ejercicio 17 =====
\begin{ejercicio}
Calcule $A \times B$, donde
\[
A=\begin{bmatrix} 2 & -1 \\ 0 & 3 \\ 4 & 5 \end{bmatrix}, \quad
B=\begin{bmatrix} 1 & 4 & 2 \\ -2 & 0 & 3 \end{bmatrix}.
\]
\end{ejercicio}
\muestraSolucion{
\[
AB=\begin{bmatrix} 4 & 8 & 1 \\ -6 & 0 & 9 \\ -6 & 16 & 23 \end{bmatrix}.
\]
}

% ===== Ejercicio 18 =====
\begin{ejercicio}
Calcule $A \times B$, donde
\[
A=\begin{bmatrix} 1 & -2 & 0 \\ 3 & 1 & 4 \end{bmatrix}, \quad
B=\begin{bmatrix} 2 & 3 \\ 0 & -1 \\ 5 & 2 \end{bmatrix}.
\]
\end{ejercicio}
\muestraSolucion{
\[
AB=\begin{bmatrix} 2 & 5 \\ 26 & 16 \end{bmatrix}.
\]
}

% ===== Ejercicio 19 =====
\begin{ejercicio}
Calcule $A \times B$, donde
\[
A=\begin{bmatrix} 1 & 0 & 2 \\ 4 & -1 & 3 \\ 2 & 2 & 1 \end{bmatrix}, \quad
B=\begin{bmatrix} 3 & 1 \\ 0 & -2 \\ 1 & 4 \end{bmatrix}.
\]
\end{ejercicio}
\muestraSolucion{
\[
AB=\begin{bmatrix} 5 & 9 \\ 15 & 18 \\ 7 & 2 \end{bmatrix}.
\]
}

% ===== Ejercicio 20 =====
\begin{ejercicio}
Calcule $A \times B$, donde
\[
A=\begin{bmatrix} 3 & 5 & 1 \\ 0 & -2 & 4 \\ 1 & 1 & 1 \end{bmatrix}, \quad
B=\begin{bmatrix} 2 & 0 & 3 \\ 1 & -1 & 2 \\ 0 & 4 & 1 \end{bmatrix}.
\]
\end{ejercicio}
\muestraSolucion{
\[
AB=\begin{bmatrix} 11 & -1 & 20 \\ -2 & 18 & 0 \\ 3 & 3 & 6 \end{bmatrix}.
\]
}